% ============================================================
% 第 23 章 直流发电机（教材第六篇，印刷页 254--258）
% 节名：第一节 直流发电机的基本电磁关系；第二节 他励直流发电机的运行特性；第三节 并励直流发电机的自励及运行特性
% ============================================================
\chapter{直流发电机}

直流发电机由原动机拖动，把机械能转换为直流电能。它已大量被整流电源取代，但其"自励建压"的物理过程仍是直流电机理论中最精彩的部分，也是理解并励电动机与励磁系统的钥匙。

\section{直流发电机的基本电磁关系}

\subsection{电动势与电压方程}
电枢旋转感应电动势 $E_a=C_e\Phi n$。作为发电机，电枢输出电流，故电枢回路有电阻压降，端电压为
\begin{equation}
U=E_a-I_aR_a ,
\end{equation}
其中 $R_a$ 为电枢回路总电阻（含电枢绕组电阻与电刷接触电阻）。

\subsection{电磁转矩}
$T=C_T\Phi I_a$。发电机中电磁转矩是\textbf{制动转矩}，与原动机拖动转矩方向相反；稳态时 $T_1=T+T_0$。

\subsection{功率关系}
\begin{equation}
P_1=P_{em}+p_{mec}+p_{Fe}+p_{ad},\qquad P_{em}=E_aI_a,\qquad P_2=UI_a=P_{em}-p_{Cu a},
\end{equation}
其中 $P_1$ 为原动机输入功率，$P_2$ 为电输出功率，$p_{mec}$ 为机械损耗，$p_{Fe}$ 为铁耗，$p_{ad}$ 为杂散损耗，$p_{Cua}=I_a^2R_a$ 为电枢铜耗（对并励机还应计入励磁损耗 $U I_f$）。

\subsection{功率流程与效率}
把上述关系画成流程：原动机输入 $P_1$ 扣除机械损耗、铁耗与杂散损耗后成为\textbf{电磁功率} $P_{em}=E_aI_a$；电磁功率扣除电枢铜耗后即为\textbf{输出电功率} $P_2=UI_a$。效率为
\begin{equation}
\eta=\frac{P_2}{P_1}\times100\%=\frac{UI_a}{UI_a+p_{Cua}+p_{mec}+p_{Fe}+p_{ad}}\times100\% .
\end{equation}
其中铁耗与机械损耗近似为\textbf{不变损耗}（与负载无关，只与转速、磁通有关），铜耗为\textbf{可变损耗}（$\propto I_a^2$）。与变压器、感应电机一样，当可变损耗等于不变损耗时效率最高。

对并励发电机，励磁功率由电枢自身供给，故输出电流 $I=I_a-I_f$，输出功率 $P_2=UI$；计算效率时既可用"输入机械功率"作分母，也可用"输出电功率 $+$ 全部损耗"作分母，两者必须一致（不要重复计入励磁损耗）。

\section{他励直流发电机的运行特性}

他励发电机的励磁由独立电源供给，$I_f$ 可独立调节，是最基本的分析对象。

\subsection{空载特性}
$n=n_N$、$I_a=0$ 时，$U_0=E_a=f(I_f)$。这就是电机的磁化曲线：起始段线性（气隙线），随后弯曲饱和。空载特性是求磁路饱和程度与电枢反应大小的基础。

\subsection{外特性}
$n=n_N$、$I_f=$ 常数时，$U=f(I_a)$。随负载电流增大，$U$ 下降，原因有二：① 电枢电阻压降 $I_aR_a$；② 电枢反应的去磁作用使 $\Phi$ 略减。外特性是一条略向下弯的曲线。

\subsection{调节特性}
$n=n_N$、$U=$ 常数时，$I_f=f(I_a)$。为维持端电压不变，负载增加时须增大励磁电流，故曲线略上翘。

\subsection{电压变化率}
从空载到额定负载，端电压的相对下降称为电压变化率：
\begin{equation}
\Delta U=\frac{U_0-U_N}{U_N}\times100\% ,
\end{equation}
他励发电机一般约 5\%～10\%。

\section{并励直流发电机的自励及运行特性}

\subsection{自励建压过程}
并励发电机的励磁绕组与电枢并联，励磁电流取自电枢端电压。它不需要外部直流电源即可建立电压，过程如下：
\begin{enumerate}
  \item 电机由原动机拖动到额定转速，电枢绕组切割主磁极的\textbf{剩磁}，感应一个不大的剩磁电动势 $E_r$（约为额定电压的 2\%～5\%）；
  \item $E_r$ 加在励磁绕组上，产生不大的励磁电流 $I_f$；
  \item 若该励磁电流产生的磁通与剩磁\textbf{同方向}，则磁通增强，$E_a$ 增大，$I_f$ 进一步增大……
  \item 如此循环，电压不断上升，直到励磁绕组的伏安特性 $U=I_fR_f$ 与电机的空载特性 $E_0=f(I_f)$ 相交于稳定点为止。
\end{enumerate}

\subsection{自励的三个条件}
\begin{enumerate}
  \item 主磁极必须\textbf{有剩磁}（否则须先用外电源充磁）；
  \item 励磁绕组与电枢的\textbf{接法必须正确}，使励磁电流产生的磁通与剩磁同向（接反则电压建不起来，甚至反转极性）；
  \item 励磁回路的\textbf{电阻必须小于临界电阻} $R_{f,cr}$。若 $R_f$ 过大，励磁特性线 $U=I_fR_f$ 与空载特性的交点过低，电压只能建立到很低的数值。
\end{enumerate}

\keyformula{R_f<R_{f,cr}\quad\text{（自励建压的必要条件）}}

\subsection{外特性}
并励发电机的外特性 $U=f(I_a)$ 比他励机\textbf{更陡}：负载增加时，除电枢电阻压降与电枢反应去磁外，端电压下降还使励磁电流减小，进一步削弱磁场。因此并励机存在\textbf{临界电流}：负载电阻减小时，电流增大到某一值后反而减小（因电压崩溃）。

\subsection{建压失败的排查}
并励发电机建不起电压是常见的故障现象，按自励三条件逐一排查：

\begin{center}
\begin{tabular}{@{}p{3.2cm}p{4.0cm}p{4.8cm}@{}}
\toprule
现象 & 可能原因 & 处理 \\
\midrule
电压表指示为零或极小 & 无剩磁 & 用外电源充磁 \\
电压建立但极性相反 & 励磁绕组接反 & 对调励磁绕组两端 \\
电压只能建立到很低值 & 励磁回路电阻过大（超过临界电阻） & 减小磁场变阻器阻值 \\
电压建立后不稳、振荡 & 转速过低或电刷接触不良 & 提高转速、研磨/更换电刷 \\
\bottomrule
\end{tabular}
\end{center}

\subsection{并励与他励的比较}
\begin{center}
\begin{tabular}{@{}p{2.2cm}p{4.0cm}p{4.4cm}@{}}
\toprule
项目 & 他励发电机 & 并励发电机 \\
\midrule
励磁电源 & 独立直流电源 & 取自电枢端电压 \\
建压 & 加励磁即有电压 & 需剩磁自励建压 \\
外特性 & 略下降 & 下降更陡，有临界电流 \\
适用 & 宽范围调压、试验 & 常用直流电源 \\
\bottomrule
\end{tabular}
\end{center}

\corepoint{自励建压这一节是"正反馈"的经典例子：剩磁 $\rightarrow$ 小电动势 $\rightarrow$ 小励磁电流 $\rightarrow$ 磁通增强 $\rightarrow$ 电动势增大 $\rightarrow$ …直到饱和与电阻压降把增长"刹住"。三个条件缺一不可：\textbf{有剩磁、接法对（助磁）、电阻小于临界值}。考试常问"并励发电机建不起电压的可能原因"，就按这三条逐一排查（另外还要检查转速是否过低、电刷接触是否良好）。}
